\bookchapter{Bab 47: Panduan Saku dan Sumber Belajar}

\chapterintro{Bab ini adalah panduan saku yang bisa dipakai secara mandiri. Saat kamu lupa suatu prinsip atau perlu mengambil keputusan dengan cepat, langsung buka bagian ini. Saran kami: perlakukan seperti “kartu saku”—ringkasan lima prinsip utama, pohon keputusan alat, pertanyaan umum, dan sumber belajar, semuanya lengkap tanpa terlewat satu hal pun.}

\booksection{Kartu saku model berpikir}

\begin{center}
\begin{tabularx}{\linewidth}{|X|X|X|}
\hline
\textbf{Prinsip} & \textbf{Satu kalimat} & \textbf{Poin praktis} \tabularnewline
\hline
Think Before Coding & Pikirkan matang dulu sebelum bicara & Jawab tiga pertanyaan: apa yang dikerjakan? seperti apa hasilnya? apa batasannya? \tabularnewline
Simplicity First & Selesaikan masalah dengan kode sesedikit mungkin & Satu fitur saja setiap kali, buang kompleksitas yang tidak perlu \tabularnewline
Surgical Changes & Ubah secara presisi tanpa melampaui batas & Tentukan secara eksplisit cakupan perubahan, tolak "sekalian saya ubahkan ini" dari AI \tabularnewline
Goal-Driven Execution & Sebut "apa yang dikerjakan", bukan "bagaimana mengerjakannya" & Beri tujuan, bukan langkah-langkahnya; biarkan AI memilih jalannya sendiri \tabularnewline
Iterate \& Verify & Siklus hasilkan-verifikasi-koreksi & Verifikasi tiga lapis: fungsi, batas, arsitektur \tabularnewline
\hline
\end{tabularx}
\end{center}

\begin{center}\small\color{gray}Kartu saku lima prinsip berpikir utama\end{center}

\booksection{Pohon keputusan pemilihan alat}

Saat kamu perlu memilih alat pemrograman AI, ikuti jalur keputusan berikut. Perhatikan: ini bukan peringkat “mana yang terbaik”, melainkan kerangka penilaian “mana yang paling cocok untukmu”.

Anda ingin membuat apa?

Nol dasar, ingin cepat merasakan pilih Builder cloud (v0/Bolt/Replit)

Pengembangan harian, butuh pelengkap kode pilih AI bawaan IDE (Cursor/Trae)

Proyek besar, butuh operasi massal pilih Agent terminal (Claude Code)

Uji ide cepat lewat prototipe pilih Builder cloud atau Trae SOLO

\begin{center}\small\color{gray}Gambar 6: Pohon keputusan pemilihan alat pemrograman AI\end{center}

\booksection{Pertanyaan umum (FAQ)}

Q1: Saya sama sekali tidak bisa pemrograman, bisakah tetap coding dengan AI?

Bisa. Alat pemrograman AI tahun 2026 sudah cukup kuat; kamu bisa mendeskripsikan kebutuhan dengan bahasa alami dan membiarkan AI menghasilkan kodenya. Tetapi ada dua saran: pertama, mulai dari alat Builder berbasis cloud (seperti v0, Bolt) yang tidak butuh konfigurasi lingkungan pengembangan; kedua, pelajari konsep dasar sambil memakainya—tidak perlu mahir sintaks, tetapi kamu harus bisa “mengerti garis besar” kode yang ditulis AI agar penerimaannya efektif.

Q2: Bahasa pemrograman apa yang sebaiknya saya pelajari?

Tergantung apa yang mau kamu kerjakan. Frontend situs web—JavaScript/TypeScript; pengolahan data atau AI—Python; backend tingkat enterprise—Java atau Go. Tetapi yang lebih penting: di era pemrograman AI, pentingnya “bahasa” sedang menurun. AI membantu menangani detail sintaks; yang perlu kamu pahami adalah “struktur program”—apa itu fungsi, apa itu perulangan, apa itu struktur data. Konsep-konsep ini berlaku lintas bahasa; begitu paham, berpindah ke bahasa apa pun jadi cepat.

Q3: Apakah kode buatan AI bermasalah dari sisi hak cipta?

Sampai Agustus 2026, ini masih menjadi persoalan hukum yang diperdebatkan. Pada umumnya, kode yang dihasilkan AI tidak dianggap memiliki pemilik hak cipta yang jelas di sebagian besar yurisdiksi. Tetapi jika dipakai untuk proyek komersial, saran kami: (1) pelajari ketentuan kekayaan intelektual dalam syarat layanan alat AI yang kamu pakai; (2) lakukan tinjauan dan modifikasi manual atas kode-kode penting agar memuat masukan kreatifmu; (3) hindari membiarkan AI menyalin langsung kode sumber terbuka yang berpotensi berhak cipta.

Q4: Apakah AI akan menggantikan para pemrogram?

Dalam jangka pendek, tidak. Yang digantikan AI bukan profesi “pemrogram”, melainkan pekerjaan “mengetik kode baris demi baris”. Persis seperti kalkulator tidak menggantikan matematikawan, melainkan membebaskan mereka dari perhitungan manual yang merepotkan. Alat pemrograman AI menggeser pemrogram dari “menulis kode” ke “memikirkan masalah”—ini justru menaikkan batas atas nilai seorang pemrogram. Daya saing di masa depan tidak terletak pada “bisa menulis berapa baris kode”, melainkan pada “bisa memecahkan masalah seberapa rumit”.

Q5: Bagaimana menilai apakah hasil kerja AI sudah benar?

Metode verifikasi tiga lapis. Lapis pertama: verifikasi fungsi—bisakah ia dijalankan? Benarkah input dan outputnya? Ini yang paling dasar. Lapis kedua: verifikasi batas—apa yang terjadi dengan input kosong, input super panjang, atau input aneh? Ini menemukan masalah yang tersembunyi. Lapis ketiga: verifikasi arsitektur—apakah struktur kode masuk akal? Adakah yang ditulis langsung (hardcoded)? Apakah mudah dikembangkan? Ini menuntut daya nilaimu secara teknis. Kalau kamu belum bisa menilai lapis ketiga, minimal penuhi dua lapis pertama—bisa dijalankan dan batasnya normal, itu sudah lebih baik daripada kebanyakan pemula.

\booksection{Sumber belajar yang direkomendasikan}

Berikut sumber belajar rekomendasi per Agustus 2026. Perlu diingat, tautan internet bisa berubah; jika ada tautan yang mati, biasanya mencari judulnya akan menemukan versi terbarunya.

\begin{center}
\begin{tabularx}{\linewidth}{|X|X|X|}
\hline
\textbf{Jenis} & \textbf{Sumber} & \textbf{Alasan direkomendasikan} \tabularnewline
\hline
Sumber primer & CLAUDE.md milik Andrej Karpathy & Sumber dari kelima prinsip; hanya 70 baris tetapi setiap katanya berharga \tabularnewline
Sumber primer & Dokumentasi resmi Anthropic (Claude Code) & Panduan praktik terbaik untuk Agent di terminal \tabularnewline
Tutorial & Koleksi Skills yang dihimpun komunitas Vibe Coding & Kumpulan trik praktik yang dirangkum para praktisi komunitas \tabularnewline
Tutorial & Dokumentasi resmi dan kanal YouTube Cursor & Panduan otoritatif memakai AI bawaan IDE \tabularnewline
Komunitas & awesome-claude-code di GitHub & Praktik terbaik dan koleksi plugin Claude Code \tabularnewline
Komunitas & Komunitas r/cursor di Reddit & Berbagi pengalaman dan diskusi pengguna Cursor dari seluruh dunia \tabularnewline
Dasar & MDN Web Docs (developer.mozilla.org) & Dokumentasi rujukan paling otoritatif untuk pengembangan frontend \tabularnewline
Dasar & Tutorial resmi Python (docs.python.org) & Titik awal resmi untuk belajar Python \tabularnewline
\hline
\end{tabularx}
\end{center}

\begin{center}\small\color{gray}Daftar sumber belajar yang direkomendasikan\end{center}

\booksection{Penutup: alat di tangan, metodologi di hati}

Kembali pada keyakinan inti buku ini: alat beriterasi setiap tiga bulan, metodologi berevolusi setiap tiga tahun. Pada hari ini, Agustus 2026, kemampuan alat pemrograman AI sudah kuat luar biasa—kamu bisa membuat situs web dengan satu kalimat, menghasilkan aplikasi full-stack dengan satu deskripsi, menyelesaikan restrukturisasi rumit dengan satu instruksi. Tetapi semakin kuat alatnya, justru semakin menonjol pentingnya metodologi—sebab kekuatan alat memperbesar setiap keputusan yang kamu ambil.

Jika kebutuhanmu jelas, AI membantumu mewujudkannya dengan efisien; jika kebutuhanmu samar, AI juga dengan efisien menghasilkan hal yang salah. Jika kamu patuh pada kesederhanaan, AI membantumu memverifikasi gagasan dengan cepat; jika kamu serakah menghendaki segalanya, AI juga membantumu menumpuk utang teknis dengan cepat. Jika kamu mengubah dengan presisi, AI membantumu mengoptimalkan kode sepresisi bedah; jika instruksimu samar, AI juga membantumu memasukkan bug dengan gencar.

Itulah sebabnya buku ini menghabiskan sebagian besar halamannya untuk metodologi, bukan operasional alat. Bukan karena alat tidak penting, melainkan karena metodologi adalah “pengali” bagi alat—semakin baik metodologinya, semakin besar nilai alat itu; jika metodologinya nol, alat sekuat apa pun tetap nol dikali bilangan apa pun.

Anda tidak lagi coding sendirian. Anda seorang komandan, AI adalah pasukan Anda. Kualitas cara berpikir Anda menentukan kekuatan tempur pasukan itu. Mulai hari ini, bersenjatailah dengan metodologi dan perbesar kemampuan Anda dengan alat—itulah seluruh rahasia coding bersama AI.

\booksection{Daftar pustaka}

\begin{enumerate}
  \item Andrej Karpathy. (2025). Vibe Coding. Social media post introducing the concept of vibe-driven programming with AI.
  \item Andrej Karpathy. (2025). CLAUDE.md - Project guidelines for AI-assisted development. GitHub repository, 70+K stars.
  \item Anthropic. (2026). Claude Code: Terminal-native AI programming. Official documentation and best practices.
  \item Cursor Team. (2026). Cursor 3: The next generation of AI-first IDE. Release notes and documentation.
  \item OpenAI. (2026). Codex CLI: Bringing AI to the terminal. Technical documentation.
  \item Model Context Protocol (MCP) Specification. (2025-2026). Open protocol for AI-tool integration.
  \item Vibe Coding Community. (2026). Community-curated skills and best practices for AI-assisted programming.
  \item Spec-Driven Development (SDD) methodology. (2026). Emerging practices in specification-first AI programming.
\end{enumerate}
